% GENERATED FILE. Edit canonical lessons, then run npm run generate:books.
% canonical-source-hash: fnv1a64:c11506dfb26f4b9f
% canonical-lessons: MW-C280-tim, MW-C280-ged, MW-C280-jhelno, MW-C280-nisano, MW-C280-sadhno

\chapter{Ball Games}
\label{ch:mw-ball-games}
\hlchaptermodality{280}

\begin{chapteropening}
\textbf{By the end of this chapter:} \emph{I can say what I do on the playing field.}

Complete the last lesson of chapter 280: I can say what I do on the playing field.
\end{chapteropening}

\section[\texorpdfstring{\mw{टीम} (ṭīm)}{टीम (ṭīm)}]{\mw{टीम} (ṭīm) — a team}

\label{lesson:MW-C280-tim}



\begin{quote}
Before the new one: say the Marwadi for a contest, then the Marwadi for wrestling.
\end{quote}

\subsection*{You'll want to know: \mw{टीम}}
\textbf{\mw{टीम}} — \emph{ṭīm} — \textquotedblleft{}a team\textquotedblright{}.

\begin{quote}
\mw{म्हारी} \mw{टीम} \mw{जीतै} \mw{है।}
\end{quote}

\begin{quote}
\emph{mhārī ṭīm jītai hai.}
\end{quote}

\begin{quote}
My team wins.
\end{quote}

A feminine noun.

\begin{cousinweb}
\textbf{\mw{टीम}} (\emph{ṭīm}) is borrowed from English \emph{team}. It is feminine: \textbf{\mw{म्हारी} \mw{टीम}}.
\end{cousinweb}

\begin{grammarlens}[title={What you've built}]
One more word for this chapter.
\end{grammarlens}

\subsection*{Your turn}
\begin{itemize}
\raggedright
  \item \emph{Say it:} \emph{ṭīm}
  \item \emph{Say it:} \emph{ṭīm}, once more
  \item \emph{Say it:} \emph{mhārī ṭīm jītai hai.} — My team wins.
  \item \emph{Recall:} say the Marwadi for a contest, then the Marwadi for wrestling, then say \emph{ṭīm} again
\end{itemize}

\begin{tcolorbox}[breakable,colback=teal!4,colframe=teal!35!black,title={Before you move on}]
What does \mw{टीम} mean? (\textquotedblleft{}A team\textquotedblright{}.) Say it once more.
\end{tcolorbox}
\section[\texorpdfstring{\mw{गेंद} (gẽd)}{गेंद (gẽd)}]{\mw{गेंद} (gẽd) — a ball}

\label{lesson:MW-C280-ged}



\begin{quote}
Before the new one: say the Marwadi for wrestling, then the Marwadi for a team.
\end{quote}

\subsection*{You'll want to know: \mw{गेंद}}
\textbf{\mw{गेंद}} — \emph{gẽd} — \textquotedblleft{}a ball\textquotedblright{}.

\begin{quote}
\mw{म्हारी} \mw{गेंद} \mw{लाल} \mw{है।}
\end{quote}

\begin{quote}
\emph{mhārī gẽd lāl hai.}
\end{quote}

\begin{quote}
My ball is red.
\end{quote}

A feminine noun.

\begin{cousinweb}
\textbf{\mw{गेंद}} (\emph{gẽd}) is inherited from Sanskrit \emph{genduka}, \textquotedblleft{}a ball for play\textquotedblright{}; CDIAL gives no Marwari form in this entry, but lists Hindi \emph{gẽd}, \textquotedblleft{}ball\textquotedblright{} (CDIAL 4248).
\end{cousinweb}

\begin{grammarlens}[title={What you've built}]
One more word for this chapter.
\end{grammarlens}

\subsection*{Your turn}
\begin{itemize}
\raggedright
  \item \emph{Say it:} \emph{gẽd}
  \item \emph{Say it:} \emph{gẽd}, once more
  \item \emph{Say it:} \emph{mhārī gẽd lāl hai.} — My ball is red.
  \item \emph{Recall:} say the Marwadi for wrestling, then the Marwadi for a team, then say \emph{gẽd} again
\end{itemize}

\begin{tcolorbox}[breakable,colback=teal!4,colframe=teal!35!black,title={Before you move on}]
What does \mw{गेंद} mean? (\textquotedblleft{}A ball\textquotedblright{}.) Say it once more.
\end{tcolorbox}
\section[\texorpdfstring{\mw{झेलणो} (jhelṇo)}{झेलणो (jhelṇo)}]{\mw{झेलणो} (jhelṇo) — to catch a ball}

\label{lesson:MW-C280-jhelno}



\begin{quote}
Before the new one: say the Marwadi for a team, then the Marwadi for a ball.
\end{quote}

\subsection*{You'll want to know: \mw{झेलणो}}
\textbf{\mw{झेलणो}} — \emph{jhelṇo} — \textquotedblleft{}to catch a ball\textquotedblright{}.

\begin{quote}
\mw{भाई} \mw{गेंद} \mw{झेलै} \mw{है।}
\end{quote}

\begin{quote}
\emph{bhāī gẽd jhelai hai.}
\end{quote}

\begin{quote}
My brother catches the ball.
\end{quote}

Past (I …): \mw{म्हैं} \mw{गेंद} \mw{झेली}, \emph{mhaĩ gẽd jhelī}.

\begin{cousinweb}
\textbf{\mw{झेलणो}} (\emph{jhelṇo}) is inherited from a reconstructed \emph{*jhēlati}, \textquotedblleft{}undergoes\textquotedblright{}; CDIAL lists Old Marwari \emph{jhelaï}, \textquotedblleft{}to undergo, suffer\textquotedblright{}, Hindi \emph{jhelnā}, \textquotedblleft{}to take\textquotedblright{}, and Marathi \emph{jhelṇẽ}, \textquotedblleft{}to catch\textquotedblright{} (CDIAL 5413). CDIAL gives the catching sense only for Marathi; the Marwadi sense here is the everyday one.
\end{cousinweb}

\begin{grammarlens}[title={What you've built}]
One more word for this chapter.
\end{grammarlens}

\subsection*{Your turn}
\begin{itemize}
\raggedright
  \item \emph{Say it:} \emph{jhelṇo}
  \item \emph{Say it:} \emph{jhelṇo}, once more
  \item \emph{Say it:} \emph{bhāī gẽd jhelai hai.} — My brother catches the ball.
  \item \emph{Recall:} say the Marwadi for a team, then the Marwadi for a ball, then say \emph{jhelṇo} again
\end{itemize}

\begin{tcolorbox}[breakable,colback=teal!4,colframe=teal!35!black,title={Before you move on}]
What does \mw{झेलणो} mean? (\textquotedblleft{}To catch a ball\textquotedblright{}.) Say it once more.
\end{tcolorbox}
\section[\texorpdfstring{\mw{निसाणो} (nisāṇo)}{निसाणो (nisāṇo)}]{\mw{निसाणो} (nisāṇo) — a target; an aim}

\label{lesson:MW-C280-nisano}



\begin{quote}
Before the new one: say the Marwadi for a ball, then the Marwadi for to catch a ball.
\end{quote}

\subsection*{You'll want to know: \mw{निसाणो}}
\textbf{\mw{निसाणो}} — \emph{nisāṇo} — \textquotedblleft{}a target; an aim\textquotedblright{}.

\begin{quote}
\mw{निसाणो} \mw{दूर} \mw{है।}
\end{quote}

\begin{quote}
\emph{nisāṇo dūr hai.}
\end{quote}

\begin{quote}
The target is far away.
\end{quote}

A masculine noun.

\begin{cousinweb}
\textbf{\mw{निसाणो}} (\emph{nisāṇo}) is usually traced to Persian \emph{nishāna}, \textquotedblleft{}a mark, a target\textquotedblright{}, built on \emph{nishān}, \textquotedblleft{}a sign\textquotedblright{}; Marwadi gives it the -\mw{ो} ending of a masculine noun.
\end{cousinweb}

\begin{grammarlens}[title={What you've built}]
One more word for this chapter.
\end{grammarlens}

\subsection*{Your turn}
\begin{itemize}
\raggedright
  \item \emph{Say it:} \emph{nisāṇo}
  \item \emph{Say it:} \emph{nisāṇo}, once more
  \item \emph{Say it:} \emph{nisāṇo dūr hai.} — The target is far away.
  \item \emph{Recall:} say the Marwadi for a ball, then the Marwadi for to catch a ball, then say \emph{nisāṇo} again
\end{itemize}

\begin{tcolorbox}[breakable,colback=teal!4,colframe=teal!35!black,title={Before you move on}]
What does \mw{निसाणो} mean? (\textquotedblleft{}A target; an aim\textquotedblright{}.) Say it once more.
\end{tcolorbox}
\section[\texorpdfstring{\mw{सांधणो} (sā̃dhṇo)}{सांधणो (sā̃dhṇo)}]{\mw{सांधणो} (sā̃dhṇo) — to take aim}

\label{lesson:MW-C280-sadhno}



\begin{quote}
Before the new one: say the Marwadi for to catch a ball, then the Marwadi for a target.
\end{quote}

\subsection*{You'll want to know: \mw{सांधणो}}
\textbf{\mw{सांधणो}} — \emph{sā̃dhṇo} — \textquotedblleft{}to take aim\textquotedblright{}.

\begin{quote}
\mw{छोरो} \mw{निसाणो} \mw{सांधै} \mw{है।}
\end{quote}

\begin{quote}
\emph{chhoro nisāṇo sā̃dhai hai.}
\end{quote}

\begin{quote}
The boy takes aim at the target.
\end{quote}

Past (I …): \mw{म्हैं} \mw{निसाणो} \mw{सांध्यो}, \emph{mhaĩ nisāṇo sā̃dhyo}.

\begin{cousinweb}
\textbf{\mw{सांधणो}} (\emph{sā̃dhṇo}) is inherited from Sanskrit \emph{samdadhāti}, \textquotedblleft{}places on, combines\textquotedblright{}; CDIAL lists an Old Marwari form of this verb glossed \textquotedblleft{}aims\textquotedblright{}, and Gujarati \emph{sā̃dhvũ}, \textquotedblleft{}to unite, attach\textquotedblright{} (CDIAL 12898).
\end{cousinweb}

\begin{grammarlens}[title={What you've built}]
One more word for this chapter.
\end{grammarlens}

\subsection*{Your turn}
\begin{itemize}
\raggedright
  \item \emph{Say it:} \emph{sā̃dhṇo}
  \item \emph{Say it:} \emph{sā̃dhṇo}, once more
  \item \emph{Say it:} \emph{chhoro nisāṇo sā̃dhai hai.} — The boy takes aim at the target.
  \item \emph{Recall:} say the Marwadi for to catch a ball, then the Marwadi for a target, then say \emph{sā̃dhṇo} again
\end{itemize}

\begin{tcolorbox}[breakable,colback=teal!4,colframe=teal!35!black,title={Before you move on}]
What does \mw{सांधणो} mean? (\textquotedblleft{}To take aim\textquotedblright{}.) Say it once more.
\end{tcolorbox}
